% Opens the second part: what the map is and how to read its entries.
% Labels of this file begin with kind:reading.
\chapter{How to read the map}
\label{chap:reading}

The \emph{map} is the rest of the Lean development of this repository: the library
\texttt{SagittaMap} (\texttt{Sagitta/Map}, namespace \texttt{Sagitta.Map}). It is where the ideas
of the arrow were tried before the chain of Part~I was written down, and it keeps what was
formalized on the way. The arrow does not import the map: nothing in Part~I depends on an entry of
this part.

\textbf{Language.} The identifiers and the comments of the map are in Spanish, as they were
written. Each entry of this part gives the statement in English next to the Lean name, and says in
which file it is.

\textbf{Kinds of entries.} An entry of the map is one of four kinds, and says which:
\begin{itemize}
  \item a result of the literature, formalized again here: it names its \emph{source} and carries a
    \emph{Fidelity} note;
  \item a definition or a statement of this project, marked as its \emph{own}, with no literature
    source;
  \item \emph{infrastructure}: a tool that other entries use (an encoding, a closure property, a
    bridge between two libraries), not a result in itself;
  \item an \emph{exploration}: the record of an attempt, kept as it was left.
\end{itemize}

\textbf{Models.} Many results of the map are not stated for Turing machines. They are stated in a
smaller model (a tree of configurations, an abstract cost model, a Boolean circuit, a finite
structure), in which the argument of the source can be followed without the machine. Each entry
says in which model its Lean statement lives, and the Fidelity note says how that differs from the
statement of the source. A result in a smaller model is not the theorem of the source for
machines, and no entry claims that it is.

\textbf{Fidelity.} A Fidelity note that says ``Not checked against the printed source'' means
exactly that: the Lean statement was written from the standard form of the result and has not been
compared with the publication. Chapter~\ref{chap:fidelity} lists the entries with a source and the
state of each comparison.

\textbf{What a green node means in this part.} The definition or the statement is formalized and
its proof is complete: the map builds with no \texttt{sorry}. As in Part~I, a named hypothesis is a
definition of type \texttt{Prop} taken as an explicit argument; a green node for it says that the
definition exists, not that the hypothesis holds. Chapter~\ref{chap:hypotheses} lists the named
hypotheses of both parts.
% The file-by-file check of axioms (#print axioms under #guard_msgs) that Part I describes is not
% in place for the map yet. When it is, say so here.
